\scanpage{82}
\mspage{4}
pris du Kimchi dans un saladier et qu'il en reste après le repas, ne pas\\
le remettre dans le pot, mais le garder pour le terminer à l'un des repas\\
suivants. Il se gardera plusieurs jours dans le saladier sans problème,\\
alors que le Kimchi dans le pot doit se garder plusieurs semaines, voire\\
plusieurs mois.

Comme matériel pour la préparation proprement dite, il faut prévoir\\
une planche pour découper les légumes, un grand couteau de cuisine, plu-\\
sieurs sceaux ou grandes cuvettes pour le lavage \unsure{et} le malaxage des légumes,\\
un ou deux grands tamis ronds en paille ou bambou tressés pour égoutter -\\
on en trouve d'importation chinoise pour vingt à trente francs. J'en ai deux,\\
et même en dehors des moments de préparation du Kimchi, ils me servent pra-\\
tiquement tous les jours, pour laver des légumes ou des céréales, comme\\
\deleted{panier} corbeille à fruits et à légumes etc. Enfin, pour les Kimchis sans\\
eau, il faut un mortier avec pilon pour écraser ensemble le sel et les aro-\\
mates.

4.\quad La préparation et la conservation du Kimchi est chose des plus simples,\\
et il n'y a aucun problème dès qu'on lui consacre un minimum d'intelligence\\
et de soin. (Celui qui n'est pas disposé à faire ce minimum d'investissement\\
fera mieux de se nourrir de conserves ...) La première et principale exigean\\
-ce est la propreté. Tous les instruments utilisés pour la préparation \unsure{et|ou}\\
l'utilisation du Kimchi doivent être propres. Plus particulièrement, au\\
moment de la préparation, le pot qu'on s'apprête à utiliser doit être lavé\\
soigneusement tant intérieurement qu'extérieurement, et séché intérieurement\\
avec un torchon propre. De même pour les sceaux et cuvettes, la planche à\\
découper etc. La table de travail doit être nette de tout objet étranger à\\
la préparation du Kimchi, et propre. Les mains et bras sont propres, et\\
tout particulièrement pour le malaxage des légumes avec les aromates (cas du\\
Kimchi sans eau). On évitera de remplir le pot à ras-bord, pour que l'assi-\\
ette-couvercle ne soit pas au contact du Kimchi, et pour éviter aussi que le\\
pot ne déborde à la suite du ``travail'' du Kimchi dans les premiers jours.

Dans les jours qui suivent la préparation du Kimchi, alors que la\\
fermentation s'amorce, il est bon de surveiller le déroulement, en remuant\\
le Kimchi à l'eau pour faire replonger les légumes montés à la surface\\
et éviter qu'ils ne se dessèchent ou ne moisissent, ou tasser (avec une\\
cuillère en bois ou une louche, ou à la main lavée au préalable) le Kimchi\\
sans eau, de sorte que le jus passe au dessus du Kimchi, dont la partie supé-\\
\unsure{rièure|rieure} restera humide. Il est bon de le faire tous les un ou deux jours.\\
U$\underset{\unsure{v}}{\overset{\unsure{-}}{\mathrm n}}$e fois que le Kimchi est bon pour la consommation, l'usage régulier du\\
Kimchi suffit à garantir en principe une bonne continuation du processus de
\mspage{81}
